%!TEX root = ../../thesis.tex
% Generated by scripts/evaluate.py -- do not edit; edit the emitter instead.
\begin{figure}[htb]
  \centering
  \begin{tikzpicture}
    \begin{axis}[thesis result plot, width=0.86\textwidth, height=6.2cm,
      xlabel={trie depth (episode step)}, ylabel={distinct prefixes reached},
      ymode=log, xmin=0, legend pos=south east, legend style={font=\scriptsize}]
    \addplot [KIT line plot A] table [col sep=comma, x=depth, y=nodes] {figures/data/tree-cjson.csv};
    \addlegendentry{cJSON}
    \addplot [KIT line plot B] table [col sep=comma, x=depth, y=nodes] {figures/data/tree-libxml2.csv};
    \addlegendentry{libxml2}
    \addplot [KIT line plot C] table [col sep=comma, x=depth, y=nodes] {figures/data/tree-picohttpparser.csv};
    \addlegendentry{picohttpparser}
    \end{axis}
  \end{tikzpicture}
  \caption[Action-trie breadth by depth]{Distinct action prefixes the run ever executed, by depth in the action-execution trie, for one representative run per target. The shape separates search that goes deep and narrow from search that goes broad and shallow; a trie that widens without deepening is the signature of the corpus formulation spending its budget on variations of the same prefix.}
  \label{fig:tree-depth}
\end{figure}
